%% ========================================================================
%% Bab 10: Pengujian dan Kualitas Aplikasi
%% ========================================================================

\chapter{Pengujian dan Kualitas Aplikasi}
\label{ch:pengujian-aplikasi}

\chapterhero{orange}{Baca anatomi test fitur Simple POS yang sudah ada, jalankan sebagai gerbang kualitas, lalu susun checklist milestone sebelum melangkah ke bab berikutnya.}{\faIcon{vial}\quad 13 feature test}{\faIcon{check-circle}\quad \cmd{php artisan test}}{\faIcon{bug}\quad checklist milestone}

Petugas quality control di pabrik tidak menunggu produk sampai ke tangan
pembeli untuk tahu ada yang cacat. Ia berdiri di ujung jalur produksi,
mengambil sampel dari setiap batch, menjalankan pemeriksaan yang sama
persis setiap kali, dan menahan batch itu di pabrik begitu ada satu saja
yang tidak lolos. Pemeriksaan itu tidak butuh insting atau tebakan;
checklist-nya sudah baku, dan hasilnya selalu berupa jawaban tegas: lolos
atau tidak. Test fitur pada Simple POS bekerja dengan peran yang sama.
Setiap kali kode berubah, entah menambah fitur baru atau memperbaiki bug
lama, menjalankan \cmd{php artisan test} memberi jawaban tegas dalam
hitungan detik: perilaku yang sudah ada sebelumnya masih berjalan
sebagaimana mestinya, atau ada yang rusak di suatu tempat yang mungkin
tidak disadari sedang disentuh.

\textbf{Yang Akan Kamu Pelajari}

\begin{enumerate}
  \item menjelaskan cakupan 13 test fitur Simple POS (login/logout, otorisasi peran, transaksi POS, alur token API) dan apa yang divalidasi tiap kelompok test;
  \item menjalankan \cmd{php artisan test} dan menafsirkan hasil 27 assertion untuk menentukan apakah sebuah perubahan kode memecahkan perilaku yang sudah ada;
  \item menyusun checklist milestone untuk memverifikasi bahwa CRUD, autentikasi, dan otorisasi Simple POS berfungsi sebelum lanjut ke bab berikutnya.
\end{enumerate}

\section{Anatomi Test Fitur Simple POS}
\label{sec:pengujian-aplikasi-1}

\begin{istilahpenting}
\begin{description}[leftmargin=0pt,style=nextline]
  \item[Feature Test] Test yang menjalankan satu alur lengkap aplikasi, dari request HTTP masuk sampai respons keluar, berbeda dari unit test yang menguji satu fungsi terisolasi.
  \item[Assertion] Satu pernyataan pemeriksaan dalam test, misalnya "status respons harus 200" atau "stok produk harus berkurang 3", yang dinyatakan gagal kalau kenyataannya tidak cocok.
  \item[RefreshDatabase] Trait PHPUnit/Pest bawaan Laravel yang mengembalikan basis data test ke kondisi bersih sebelum setiap test dijalankan, mencegah satu test memengaruhi hasil test lain lewat data sisa.
  \item[Regresi] Perilaku lama yang sebelumnya berjalan benar, tapi rusak akibat perubahan kode baru, biasanya tidak disengaja dan tidak disadari sampai test gagal.
\end{description}
\end{istilahpenting}

13 test fitur Simple POS terbagi ke empat kelompok, masing-masing
menguji satu potongan perilaku yang sudah dibahas di bab-bab
sebelumnya: autentikasi (login berhasil, login gagal, logout), otorisasi
berbasis peran (admin bisa mengakses halaman produk, kasir diblokir),
transaksi POS (stok berkurang sesuai qty, total dihitung benar dari
harga produk), dan alur token API (login API mengembalikan token, token
yang salah ditolak).

\begin{lstlisting}[language=php, title=\lstfile{tests/Feature/ProductAuthorizationTest.php}]
public function test_kasir_cannot_access_admin_pages(): void
{
    $kasir = User::factory()->create(['role' => 'kasir']);

    $response = $this->actingAs($kasir)->get('/products');

    $response->assertForbidden();
}
\end{lstlisting}

\phpclass{actingAs(\$kasir)} mensimulasikan request seolah dikirim oleh
pengguna \phpclass{\$kasir} yang sudah login, tanpa benar-benar mengisi
form login; ini yang membuat feature test bisa menguji middleware
\phpclass{role:admin} dari Bab~7 secara langsung.
\phpclass{assertForbidden()} adalah satu assertion, memeriksa bahwa
respons benar-benar berstatus 403. \phpclass{RefreshDatabase} yang
dipasang di kelas test memastikan \phpclass{\$kasir} yang baru dibuat
di test ini tidak bocor ke test lain yang berjalan setelahnya.

\section{Menjalankan dan Menafsirkan Hasil Test}
\label{sec:pengujian-aplikasi-2}

\begin{lstlisting}[language=bash]
php artisan test

   PASS  Tests\Feature\AuthenticationTest
  [x] user can login with valid credentials
  [x] user cannot login with invalid credentials

   PASS  Tests\Feature\TransactionTest
  [x] stock decrements after transaction

  Tests:    13 passed (27 assertions)
  Duration: 1.42s
\end{lstlisting}

Baris \texttt{Tests: 13 passed (27 assertions)} adalah ringkasan yang
paling penting dibaca lebih dulu: 13 adalah jumlah fungsi test yang
dijalankan, 27 adalah jumlah pernyataan pemeriksaan individual di dalam
seluruh test itu (satu fungsi test bisa berisi lebih dari satu
assertion, seperti pada \phpclass{test\_kasir\_cannot\_access\_admin\_pages}
di atas kalau ditambah pemeriksaan pesan error). Semua tanda centang
hijau berarti seluruh perilaku yang diuji masih berjalan seperti yang
diharapkan.

\begin{warningbox}
Test yang gagal setelah sebuah perubahan kode bukan gangguan yang harus
dilewati begitu saja; ia adalah tepat sinyal yang seharusnya dicari:
regresi. Kalau
\phpclass{test\_stock\_decrements\_after\_transaction} tiba-tiba gagal
setelah controller transaksi diubah, itu berarti perubahan itu merusak
perilaku pengurangan stok yang sebelumnya benar, terlepas dari apakah
perubahan itu ditujukan untuk fitur lain sama sekali.
\end{warningbox}

\section{Praktik: Checklist Milestone Sebelum Lanjut}
\label{sec:pengujian-aplikasi-3}

13 test fitur ini menutup fitur-fitur inti sampai Bab~9, tapi tidak
seluruh perilaku Simple POS punya test otomatis. Checklist manual
mengisi celah itu, terutama untuk hal-hal yang lebih mudah diperiksa
lewat browser daripada ditulis sebagai assertion. Langkah-langkah
berikut menjalankan seluruh test yang sudah ada, sengaja merusak salah
satu untuk melihat kegagalannya, lalu menuntaskan checklist manual.

\begin{enumerate}
  \item Lihat lokasi test yang sudah ada:
    \texttt{tests/Feature\allowbreak/AuthenticationTest\allowbreak.php},
    \texttt{tests/Feature\allowbreak/TransactionTest\allowbreak.php}, dan berkas
    serupa lainnya di \texttt{tests/Feature/}. Setiap berkas menguji satu
    kelompok perilaku, sesuai anatomi pada bagian sebelumnya.
  \item Jalankan seluruh suite untuk memastikan baseline masih hijau:
    \begin{lstlisting}[language=bash]
php artisan test
    \end{lstlisting}
    \textbf{Yang diharapkan}: keluaran seperti pada bagian sebelumnya,
    \texttt{13 passed (27 assertions)}.
  \item Buka \texttt{app/Http/Controllers/TransactionController.php},
    ubah sengaja baris pengurangan stok (misalnya ganti
    \texttt{\$product->decrement('stock', \$qty)} menjadi
    \texttt{\$product->increment('stock', \$qty)}, kebalikannya).
  \item Jalankan ulang \artisan{php artisan test}. \textbf{Yang
    diharapkan}: baris \texttt{Tests: 13 passed} berubah menjadi
    memuat setidaknya \texttt{1 failed}, dengan nama
    \phpclass{test\_stock\_decrements\_after\_transaction} muncul
    ditandai gagal, disertai pesan yang membandingkan nilai stok yang
    diharapkan dengan nilai yang benar-benar didapat.
  \item Kembalikan baris yang diubah di langkah 3 ke bentuk semula
    (\texttt{decrement}), jalankan \artisan{php artisan test} sekali
    lagi, dan pastikan kembali ke \texttt{13 passed}.
  \item Tuntaskan checklist manual berikut satu per satu, langsung di
    browser:
    \begin{enumerate}
      \item Buat, ubah, lalu hapus satu produk lewat \route{/products}.
        \textbf{Yang diharapkan}: ketiga aksi berhasil tanpa galat.
      \item Login sebagai admin, logout, login sebagai kasir, logout
        lagi. \textbf{Yang diharapkan}: setelah logout, membuka
        halaman yang butuh login (mis. \route{/dashboard}) langsung
        mengarahkan kembali ke form login.
      \item Login sebagai kasir, buka \route{/products}.
        \textbf{Yang diharapkan}: halaman 403, sesuai Bab~7.
      \item Catat satu transaksi lewat \route{/pos}, bandingkan stok
        produk sebelum dan sesudah di \route{/products}.
        \textbf{Yang diharapkan}: stok berkurang persis sebanyak qty
        yang dibeli.
    \end{enumerate}
  \item \textbf{Kalau test di langkah 4 tidak menunjukkan kegagalan
    sama sekali} setelah perubahan sengaja di langkah 3, periksa apakah
    baris yang diubah benar-benar baris yang dieksekusi saat transaksi
    disimpan (bukan baris pada method lain yang serupa, mis. method
    pembatalan transaksi).
\end{enumerate}

\begin{tipbox}
Jalankan \cmd{php artisan test} sebelum memulai perubahan signifikan,
bukan hanya sesudahnya. Baseline "semua hijau" sebelum perubahan
membuat kegagalan test setelahnya tidak ambigu; kalau test yang sama
sudah gagal sebelum perubahan apa pun dimulai, kegagalan itu bukan
akibat dari perubahan yang baru saja ditulis.
\end{tipbox}

\branchref{chapter-10}

\section{Rangkuman}
\label{sec:pengujian-aplikasi-rangkuman}

\begin{itemize}
  \item 13 test fitur Simple POS terbagi ke empat kelompok
    (autentikasi, otorisasi peran, transaksi POS, alur token API),
    masing-masing memakai \phpclass{actingAs()} untuk mensimulasikan
    pengguna login dan \phpclass{RefreshDatabase} untuk menjaga
    isolasi antar-test.
  \item \cmd{php artisan test} melaporkan 13 test dan 27 assertion;
    test yang gagal setelah perubahan kode adalah sinyal regresi,
    bukan gangguan yang boleh diabaikan.
  \item Checklist milestone manual, seperti verifikasi CRUD produk dan
    dekremen stok lewat browser langsung, melengkapi cakupan test
    otomatis sebelum melangkah ke bab berikutnya.
\end{itemize}

\section{Latihan}

\begin{enumerate}
  \item Jalankan \cmd{php artisan test} dan baca hasilnya baris per
    baris, catat kelompok test mana yang menguji perilaku dari Bab~7.
  \item Ubah sengaja satu baris logika pada perhitungan total transaksi
    (Bab~6) sampai salah, jalankan test lagi, dan amati test mana yang
    gagal serta pesan error apa yang ditampilkan.
  \item Tulis checklist milestone versi sendiri untuk fitur yang belum
    ada test otomatisnya di Simple POS, minimal tiga butir.
  \item \textbf{Self-check:} tanpa membuka ulang bab ini, coba jelaskan
    dengan kata-katamu sendiri: (a) perbedaan feature test dan unit
    test, (b) mengapa \phpclass{RefreshDatabase} dibutuhkan, dan (c)
    mengapa test yang gagal setelah perubahan kode disebut sinyal
    regresi. Cocokkan jawabanmu dengan isi bab ini.
\end{enumerate}

\begin{challengebox}
Tulis satu test fitur baru untuk endpoint laporan penjualan pada
Bab~8, memverifikasi bahwa filter rentang tanggal benar-benar
mengecualikan transaksi di luar rentang itu.
\end{challengebox}

\section{Referensi}

Bab ini mengacu pada dokumentasi resmi Laravel \cite{laraveldocs} untuk
feature testing dan \phpclass{RefreshDatabase}.
